% year: תש"פ
% type: מבחן
% semester: א
% moed: ג
% part: א
% number: 1
% points: 20
% categories: func-analysis, derivatives, func-limits
% source: exams/מבחנים/תשפ/מועד מיוחד תשפ.pdf

\begin{question}
חקרו את הפונקציה $y=\frac{2x-1}{(x-1)^2}$ לפי הסעיפים הבאים:
\begin{enumerate}
  \item תחום ההגדרה של הפונקציה, נקודות חיתוך עם צירים, זוגיות, אי-זוגיות.
  \item רציפות, נקודות אי-רציפות.
  \item אסימפטוטות
  \item תחום העלייה וירידה, נקודות קיצון.
  \item תחומי קמירות וקעירות, נקודות פיתול.
  \item תיאור גרפי.
\end{enumerate}
\end{question}

\begin{hint}
בגזירה לפי כלל המנה כדאי לצמצם גורם $(x-1)$ אחד מהמונה והמכנה לפני שמפשטים.
\end{hint}

\begin{solution}
נסמן $f(x)=\frac{2x-1}{(x-1)^2}$.
\begin{enumerate}
  \item \textbf{תחום הגדרה:} $D=\{x\in\R: x\neq1\}$.

  \textbf{חיתוך עם הצירים:} $f(0)=\frac{-1}{1}=-1$, כלומר $(0,-1)$ על ציר $y$. $f(x)=0\iff 2x-1=0\iff x=\frac12$, כלומר $(\frac12,0)$ על ציר $x$.

  \textbf{זוגיות:} תחום ההגדרה אינו סימטרי ביחס ל-$0$ ($-1\in D$ אבל $1\notin D$), ולכן הפונקציה אינה זוגית ואינה אי-זוגית.

  \item $f$ מנה של פולינומים, ולכן רציפה בכל תחום הגדרתה. ב-$x=1$: המונה שואף ל-$1>0$ והמכנה $(x-1)^2\to0^+$, ולכן $\lim_{x\to1^\pm}f(x)=+\infty$ --- נקודת אי-רציפות מסוג שני (אינסופית).

  \item \textbf{אנכית:} $x=1$ (משני הצדדים $f\to+\infty$).
  \textbf{אופקית:} $\lim_{x\to\pm\infty}\frac{2x-1}{(x-1)^2}=\lim_{x\to\pm\infty}\frac{\frac2x-\frac1{x^2}}{(1-\frac1x)^2}=0$, ולכן $y=0$ אסימפטוטה אופקית ב-$\pm\infty$ (ולכן אין אסימפטוטה משופעת).

  \item לפי כלל המנה:
  \[ f'(x)=\frac{2(x-1)^2-(2x-1)\cdot2(x-1)}{(x-1)^4}=\frac{2(x-1)-2(2x-1)}{(x-1)^3}=\frac{-2x}{(x-1)^3}. \]
  $f'(x)=0\iff x=0$. סימנים:
  \begin{itemize}
    \item $x<0$: מונה $-2x>0$, מכנה $<0$, ולכן $f'<0$ --- יורדת.
    \item $0<x<1$: מונה $<0$, מכנה $<0$, ולכן $f'>0$ --- עולה.
    \item $x>1$: מונה $<0$, מכנה $>0$, ולכן $f'<0$ --- יורדת.
  \end{itemize}
  לכן $f$ יורדת ב-$(-\infty,0)$ וב-$(1,\infty)$ ועולה ב-$(0,1)$. ב-$x=0$ הנגזרת עוברת משלילית לחיובית: \textbf{מינימום מקומי} $(0,-1)$. (זה גם מינימום מוחלט: עבור $x<\frac12$ הפונקציה יורדת ל-$-1$ ואז עולה, ועבור $x\ge\frac12$, $f\ge0$.)

  \item נגזור את $f'(x)=-2x(x-1)^{-3}$:
  \[ f''(x)=-2(x-1)^{-3}+6x(x-1)^{-4}=\frac{-2(x-1)+6x}{(x-1)^4}=\frac{2(2x+1)}{(x-1)^4}. \]
  המכנה חיובי לכל $x\neq1$, ולכן הסימן הוא של $2x+1$:
  $f''<0$ עבור $x<-\frac12$ --- $f$ קעורה ($\cap$) ב-$(-\infty,-\frac12)$;
  $f''>0$ עבור $-\frac12<x<1$ ו-$x>1$ --- $f$ קמורה ($\cup$) ב-$(-\frac12,1)$ וב-$(1,\infty)$.
  ב-$x=-\frac12$ הסימן מתחלף והפונקציה רציפה, ולכן זו \textbf{נקודת פיתול}: $f(-\frac12)=\frac{-2}{9/4}=-\frac89$, כלומר $(-\frac12,-\frac89)$.

  \item \textbf{תיאור גרפי:} משמאל הגרף מתחיל מתחת לציר $x$, קרוב לאסימפטוטה $y=0$ כש-$x\to-\infty$, יורד (קעור) עד נקודת הפיתול $(-\frac12,-\frac89)$, ממשיך לרדת (קמור) עד המינימום $(0,-1)$, עולה וחותך את ציר $x$ ב-$(\frac12,0)$ ושואף ל-$+\infty$ כש-$x\to1^-$. מימין ל-$x=1$ הוא יורד מ-$+\infty$ ומתקרב מלמעלה ל-$y=0$ כש-$x\to\infty$ (קמור).
\end{enumerate}
\end{solution}
